\documentclass[12pt,a4paper]{article}

% Essential packages
\usepackage[utf8]{inputenc}
\usepackage[english]{babel}
\usepackage{amsmath,amsfonts,amssymb}
\usepackage{graphicx}
\usepackage{float}
\usepackage{booktabs}
\usepackage{array}
\usepackage{multirow}
\usepackage{geometry}
\usepackage{fancyhdr}
\usepackage{parskip}
\usepackage{url}
\usepackage{hyperref}
\usepackage{listings}
\usepackage{xcolor}

% Page setup
\geometry{margin=1in}
\pagestyle{fancy}
\fancyhf{}
\rhead{\thepage}
\lhead{Spam Filter Case Study}

% Code listing setup
\definecolor{codegreen}{rgb}{0,0.6,0}
\definecolor{codegray}{rgb}{0.5,0.5,0.5}
\definecolor{codepurple}{rgb}{0.58,0,0.82}
\definecolor{backcolour}{rgb}{0.95,0.95,0.92}

\lstdefinestyle{mystyle}{
    backgroundcolor=\color{backcolour},
    commentstyle=\color{codegreen},
    keywordstyle=\color{magenta},
    numberstyle=\tiny\color{codegray},
    stringstyle=\color{codepurple},
    basicstyle=\ttfamily\footnotesize,
    breaklines=true,
    captionpos=b,
    numbers=left,
    numbersep=5pt
}

\lstset{style=mystyle}

% Title
\title{\textbf{Task 3: Development of a Spam Filter}\\
       \large A CRISP-DM Implementation for Customer Service Channels}
\author{Data Science Project}
\date{March 2026}

\begin{document}

\maketitle

\begin{abstract}
This case study presents a comprehensive machine learning-based spam filter for short message classification in customer service channels. Following CRISP-DM methodology, we developed a Multinomial Naive Bayes classifier with TF-IDF vectorization achieving 98-99\% accuracy across three risk-level configurations. The system implements adjustable decision thresholds (low-risk: 0.30, medium-risk: 0.50, high-risk: 0.70) to accommodate different business contexts and customer service scenarios. This report addresses all critical aspects: problem formulation from a business perspective, comprehensive data quality assessment, rigorous model development and evaluation, detailed error analysis with root cause identification, and an operational deployment strategy for service team integration. Key results show the medium-risk model achieves 99\% accuracy with 99\% precision while maintaining only 0.1\% false positive rate, ensuring customer experience is maintained.
\end{abstract}

\newpage
\tableofcontents
\newpage

% ============================================================================
% SECTION 1: PROBLEM CAPTURE
% ============================================================================

\section{Problem Capture: From Business Challenge to Data Science Use Case}

\subsection{Business Context and Real-World Challenge}

Short message communication (SMS, chat systems, social media feedback) has become ubiquitous in modern customer interactions. Organizations increasingly rely on these channels for customer feedback, support requests, and product inquiries. However, the open and public nature of these channels creates significant operational challenges:

\begin{itemize}
    \item \textbf{High Volume of Spam}: Unwanted messages contaminate legitimate feedback
    \item \textbf{Resource Waste}: Service teams spend significant time manually filtering messages
    \item \textbf{False Positives Impact}: Blocking legitimate customers damages relationships
    \item \textbf{Security Risk}: Phishing and malicious messages pose business risks
    \item \textbf{Operational Inefficiency}: Manual classification is time-consuming and error-prone
    \item \textbf{Variable Quality}: Customer messages range from professional to informal
\end{itemize}

\subsection{Service Team Requirements}

The service team identified critical requirements for a viable spam filtering solution:

\begin{enumerate}
    \item \textbf{High Accuracy}: Reliably distinguish spam from legitimate customer messages
    \item \textbf{Adjustable Risk Levels}: Different channels require different tolerance for false positives
    \item \textbf{Low False Positive Rate}: Blocking legitimate customers must be minimized
    \item \textbf{Operational Integration}: Solution must integrate into existing daily workflows
    \item \textbf{Transparency}: Service teams must understand and trust model decisions
    \item \textbf{Error Handling Strategy}: Clear processes for handling misclassified messages
\end{enumerate}

\subsection{Translation to Data Science Use Case}

This business problem naturally translates to a \textbf{supervised binary classification task}: Given short message text as input, predict whether the message is spam (1) or legitimate (0). The challenge requires balancing multiple objectives:

\begin{itemize}
    \item Maximizing spam detection (high recall)
    \item Minimizing false positives to protect customer relationships
    \item Accommodating multiple business contexts through risk-level adjustment
    \item Providing confidence scores for human-in-the-loop decision making
\end{itemize}

\textbf{Why CRISP-DM?} The CRISP-DM (Cross-Industry Standard Process for Data Mining) methodology provides a structured framework ensuring business requirements drive technical decisions. The six-phase approach—Business Understanding, Data Understanding, Data Preparation, Modeling, Evaluation, and Deployment—ensures we maintain alignment between technical implementation and business objectives.

\section{Why This Approach?}

Binary classification with Naive Bayes is appropriate because:

\begin{itemize}
    \item \textbf{Well-Defined Problem}: Clear distinction exists between spam and legitimate messages
    \item \textbf{Labeled Data Available}: Historical message collections with known labels exist
    \item \textbf{Proven Solution Domain}: Text classification is well-established with proven methods
    \item \textbf{Scalability}: Once trained, the model processes thousands of messages in seconds
    \item \textbf{Risk-Level Flexibility}: Threshold adjustment accommodates different business contexts without retraining
    \item \textbf{Explainability}: Probability outputs enable confidence-based decision workflows
\end{itemize}

% ============================================================================
% SECTION 2: CONCEPT AND TECHNIQUES
% ============================================================================

\section{Concept: Data Science Methods and Techniques Applied}

\subsection{Exploratory Data Analysis and Data Quality Assessment}

Before modeling, understanding data is foundational. EDA prevents blind application of algorithms to poorly understood datasets and identifies potential issues:

\textbf{Key EDA Components}:
\begin{itemize}
    \item \textbf{Dataset Composition}: Examine total messages, class distribution, data types
    \item \textbf{Class Balance Analysis}: Identify imbalanced targets requiring stratified evaluation
    \item \textbf{Descriptive Statistics}: Calculate message length, word count, character distributions
    \item \textbf{Quality Assessment}: Identify missing values, duplicates, outliers, anomalies
    \item \textbf{Vocabulary Analysis}: Identify discriminative terms and patterns
    \item \textbf{Visualization}: Create informative charts for stakeholder communication
\end{itemize}

\subsection{Text Feature Engineering: TF-IDF Vectorization}

Machine learning algorithms require numerical input. TF-IDF (Term Frequency-Inverse Document Frequency) converts text to numerical features:

\begin{equation}
\text{TF-IDF}(t,d) = \text{TF}(t,d) \times \log\left(\frac{N}{\text{document count}(t)}\right)
\end{equation}

Where TF is the frequency of term $t$ in document $d$, and IDF is the inverse document frequency.

\textbf{Why TF-IDF for this task}:
\begin{itemize}
    \item Captures word importance: Common words get lower weight, distinctive words higher weight
    \item Efficient for sparse text: Creates compact representations suitable for short messages
    \item Avoids length bias: Normalizes for message length differences
    \item Proven effectiveness: Well-established for text classification tasks
    \item Interpretability: Word importance is directly understandable by stakeholders
\end{itemize}

\textbf{Configuration Details}:
\begin{lstlisting}[language=Python]
TfidfVectorizer(
    max_features=3000,      # Limit vocabulary to top 3000 terms
    min_df=2,               # Ignore words in <2 documents
    max_df=0.8,             # Ignore words in >80% of documents
    ngram_range=(1, 2)      # Capture unigrams and bigrams
)
\end{lstlisting}

\subsection{Text Preprocessing Strategy}

Raw text contains noise that interferes with classification. Preprocessing reduces noise while preserving meaningful patterns:

\begin{enumerate}
    \item \textbf{Lowercasing}: Convert all text to lowercase ("Hello" and "hello" are equivalent)
    \item \textbf{URL Masking}: Replace URLs with [URL] token (preserves URL presence, removes variation)
    \item \textbf{Phone Number Masking}: Replace phone numbers with [PHONE] token
    \item \textbf{Punctuation Normalization}: Reduce excessive punctuation ("!!!" becomes "!!", "???" becomes "??")
    \item \textbf{Whitespace Cleaning}: Remove extra spaces and formatting
    \item \textbf{Selective Preservation}: Keep meaningful indicators like urgency words
\end{enumerate}

\subsection{Algorithm Selection: Multinomial Naive Bayes}

Naive Bayes applies Bayes' theorem with a conditional independence assumption among features:

\begin{equation}
P(y|x_1, ..., x_n) = \frac{P(x_1|y) \cdots P(x_n|y) P(y)}{P(x_1, ..., x_n)}
\end{equation}

Despite the "naive" assumption being violated in real text (words are not independent), Naive Bayes performs exceptionally well on text classification.

\textbf{Advantages for Text Classification}:
\begin{itemize}
    \item \textbf{Strong Performance}: Excellent results on short text despite simplicity
    \item \textbf{Computational Efficiency}: Very fast training and prediction
    \item \textbf{Probabilistic Output}: Provides probability estimates essential for confidence scoring
    \item \textbf{Sparse Data Friendly}: Handles high-dimensional sparse TF-IDF features well
    \item \textbf{Production Ready}: Easy to deploy and maintain
    \item \textbf{Baseline Performance}: Outperforms simpler rule-based approaches
\end{itemize}

\subsection{Risk-Level Implementation: Business-Oriented Model Configuration}

Rather than training separate models, we adjust the classification decision threshold to optimize for different business contexts. This elegant approach maintains a single model while accommodating multiple use cases.

\textbf{Decision Threshold Configuration}:
\begin{equation}
\text{Classification} = \begin{cases}
\text{SPAM} & \text{if } P(\text{spam}|x) \geq \text{threshold} \\
\text{LEGITIMATE} & \text{otherwise}
\end{cases}
\end{equation}

\begin{table}[H]
\centering
\begin{tabular}{llll}
\toprule
\textbf{Risk Level} & \textbf{Threshold} & \textbf{Use Case} & \textbf{Philosophy} \\
\midrule
Low-Risk & 0.30 & Premium/VIP customers & Very conservative \\
Medium-Risk & 0.50 & General customer service & Balanced \\
High-Risk & 0.70 & Public feedback channels & Permissive \\
\bottomrule
\end{tabular}
\caption{Risk-Level Configuration and Business Context}
\end{table}

\textbf{Business Impact}:
\begin{itemize}
    \item \textbf{Low Threshold}: Lower spam pass-through, higher false positives (conservative)
    \item \textbf{Medium Threshold}: Balanced protection and customer experience (recommended)
    \item \textbf{High Threshold}: Lower false positives, more spam pass-through (permissive)
\end{itemize}

\subsection{Evaluation Metrics and Business Interpretation}

Classification performance is measured with multiple metrics, each providing different insight:

\begin{table}[H]
\centering
\begin{tabular}{ll}
\toprule
\textbf{Metric} & \textbf{Business Meaning} \\
\midrule
Accuracy & Overall correctness: (TP+TN)/(TP+TN+FP+FN) \\
Precision & Of blocked messages, how many are truly spam: TP/(TP+FP) \\
Recall & Of all real spam, what percentage is caught: TP/(TP+FN) \\
F1-Score & Harmonic mean of precision and recall \\
False Positive Rate & Risk of blocking customers: FP/(FP+TN) \\
\bottomrule
\end{tabular}
\caption{Classification Metrics and Business Interpretation}
\end{table}

Accuracy alone is insufficient for imbalanced datasets. A model predicting only "legitimate" would achieve 86.6\% accuracy on this dataset (13.4\% spam) but have zero spam detection.

\subsection{Overfitting Prevention Strategies}

We prevent overfitting—learning training data noise rather than generalizable patterns—through multiple approaches:

\begin{enumerate}
    \item \textbf{Stratified Train-Test Split}: Maintains class distribution (13.4\% spam in both sets)
    \item \textbf{Feature Regularization}:
    \begin{itemize}
        \item max\_features=3000: Prevents learning from rare words
        \item min\_df=2: Ignores single-occurrence words (likely typos)
        \item max\_df=0.8: Ignores overly common words (low discriminative power)
    \end{itemize}
    \item \textbf{Model-Level Smoothing}: Laplace smoothing (α=0.1) prevents zero probabilities
    \item \textbf{Simple Model}: Naive Bayes' simplicity provides regularization
\end{enumerate}

% ============================================================================
% SECTION 3: ANALYSIS
% ============================================================================

\section{Analysis: Data Understanding and Model Evaluation}

\subsection{Dataset Composition and Quality Assessment}

\begin{table}[H]
\centering
\begin{tabular}{lr}
\toprule
\textbf{Metric} & \textbf{Value} \\
\midrule
Total messages & 5,572 \\
Spam messages & 747 (13.4\%) \\
Legitimate messages & 4,825 (86.6\%) \\
Training set & 4,457 (80\%) \\
Test set & 1,115 (20\%) \\
Class imbalance ratio & 6.46:1 \\
\bottomrule
\end{tabular}
\caption{Dataset Composition and Distribution}
\end{table}

\textbf{Dataset Characteristics}:

The dataset exhibits moderate class imbalance. Spam is underrepresented at 13.4%, which is realistic for customer service channels but requires careful evaluation strategies (stratified split, appropriate metrics).

\textbf{Message Length Analysis}:
\begin{itemize}
    \item Spam messages average 1,317 characters (266 words)
    \item Legitimate messages average 1,632 characters (346 words)
    \item Spam is approximately 0.81x length of legitimate messages
    \item This pattern suggests spam tends to be more concise
\end{itemize}

\textbf{Vocabulary Patterns}:
\begin{itemize}
    \item \textbf{Spam indicators}: "URGENT", "NOW", "FREE", "CLICK", "CONGRATULATIONS", "LIMITED"
    \item \textbf{Legitimate indicators}: "REMINDER", "APPOINTMENT", "ORDER", "THANK", "HELP", "SUPPORT"
    \item Distinct linguistic patterns enable effective classification
\end{itemize}

\subsection{Classification Results Across Risk Levels}

\begin{table}[H]
\centering
\begin{tabular}{@{}lcccc@{}}
\toprule
\textbf{Model} & \textbf{Accuracy} & \textbf{Precision} & \textbf{Recall} & \textbf{FP Rate} \\
\midrule
Low-Risk & 98\% & 96\% & 92\% & 0.5\% \\
Medium-Risk & 99\% & 99\% & 90\% & 0.1\% \\
High-Risk & 98\% & 99\% & 86\% & 0.1\% \\
\bottomrule
\end{tabular}
\caption{Model Performance Across Risk Levels}
\end{table}

\textbf{Medium-Risk Model (Recommended)}: Achieves 99\% accuracy with 99\% precision. This optimal configuration catches 90\% of spam (1,206 of 1,340 daily spam) while blocking only ~1 legitimate message per 10,000. This balance protects the channel while maintaining customer experience.

\textbf{Risk Level Trade-offs}:
\begin{itemize}
    \item \textbf{Low-Risk}: 96\% precision but 92\% recall (catches more spam, blocks more customers)
    \item \textbf{Medium-Risk}: 99\% precision and 90\% recall (optimal balance)
    \item \textbf{High-Risk}: 99\% precision but 86\% recall (protects customers, misses some spam)
\end{itemize}

\subsection{Detailed Error Analysis}

\begin{table}[H]
\centering
\begin{tabular}{lrr}
\toprule
 & \multicolumn{2}{c}{\textbf{Predicted}} \\
\cmidrule{2-3}
\textbf{Actual} & Legitimate & Spam \\
\midrule
Legitimate & 965 & 1 (FP) \\
Spam & 15 (FN) & 134 \\
\bottomrule
\end{tabular}
\caption{Confusion Matrix - Medium-Risk Model on 1,115 Test Messages}
\end{table}

\subsubsection{False Positive Analysis (Customer Impact)}

Only 1 legitimate message blocked out of 966 (0.1\% false positive rate). Root cause analysis identifies:

\begin{itemize}
    \item \textbf{Urgent Language}: Legitimate customers using "URGENT" or "IMMEDIATE" for real issues
    \item \textbf{Financial Discussion}: Customers complaining about billing or pricing
    \item \textbf{All-Caps Emphasis}: Frustrated customers using capitals for emphasis, not deception
    \item \textbf{Time Sensitivity}: Legitimate time-critical messages
\end{itemize}

\textbf{Mitigation Strategy}: Service teams manually review messages with low confidence (< 60\%). These blocked messages can be quickly verified and released, maintaining customer relationships.

\subsubsection{False Negative Analysis (Security Impact)}

15 spam messages passed through out of 149 spam in test set (90\% detection rate). Root causes:

\begin{itemize}
    \item \textbf{Obfuscation}: Spammers intentionally misspell words ("B0nus", "FRE3")
    \item \textbf{Sophisticated Phishing}: Well-crafted messages mimicking legitimate services
    \item \textbf{Context Blindness}: Model cannot assess sender identity or historical patterns
    \item \textbf{Evolving Tactics}: New spam patterns not represented in training data
\end{itemize}

\textbf{Mitigation Strategy}:
\begin{itemize}
    \item User reporting mechanism: Allow customers to report spam
    \item Quarterly retraining: Incorporate new spam patterns monthly
    \item Dynamic thresholds: Adjust for emerging spam operations
    \item Monitoring alerts: Flag unusual message patterns
\end{itemize}

\section{Conclusions and Operational Deployment}

\subsection{Key Results Summary}

The spam filter successfully addresses all business requirements:

\begin{itemize}
    \item \textbf{Classification Accuracy}: 98-99\% accuracy achieving high performance
    \item \textbf{Precision}: 96-99\% across risk levels ensuring spam identification confidence
    \item \textbf{Rapid Processing}: <1 millisecond per message enabling real-time classification
    \item \textbf{Business Impact}: For 10,000 messages daily:
    \begin{itemize}
        \item Medium-risk catches 1,206 spam (90\%)
        \item Blocks ~1 legitimate message (0.1\% false positive)
        \item Maintains 98.9\% legitimate customer reach
    \end{itemize}
    \item \textbf{Risk Adjustment}: Three configurations accommodate different business contexts
\end{itemize}

\subsection{Project Organization: CRISP-DM Aligned Repository Structure}

\begin{lstlisting}[language=bash, basicstyle=\ttfamily\small]
spam_filter/
├── src/                      # Source code
│   ├── spam_filter.py        # Classification engine
│   ├── gui_interface.py      # Service team GUI
│   └── report_generator.py   # Report generation
├── data/
│   ├── raw/                  # Original CSV (5,572 messages)
│   └── processed/            # training_data.json (metrics)
├── models/                   # 3 trained models
├── reports/                  # Documentation
├── figures/                  # Visualizations
└── README.md                 # Project guide
\end{lstlisting}

\subsection{Service Team Integration and Daily Workflow}

\textbf{Operational Process}:
\begin{enumerate}
    \item \textbf{High-Confidence (>80\%)}: Auto-process with monitoring (no review)
    \item \textbf{Medium-Confidence (60-80\%)}: Sample review (5-10\% manual check)
    \item \textbf{Low-Confidence (<60\%)}: Priority manual review before routing
\end{enumerate}

\textbf{Continuous Improvement Process}:
\begin{itemize}
    \item \textbf{Daily}: Monitor false positive rate, alert on anomalies
    \item \textbf{Weekly}: Review manual review queue, analyze patterns
    \item \textbf{Monthly}: Generate performance reports for management
    \item \textbf{Quarterly}: Retrain model with accumulated misclassified messages
    \item \textbf{Annual}: Evaluate new algorithms and approaches
\end{itemize}

\subsection{How the Model Integrates into Daily Work}

Each message receives: (1) Binary prediction (spam/legitimate), (2) Probability score, (3) Confidence category, (4) Business recommendation. Service teams monitor messages through an intuitive GUI, categorized by confidence level. High-confidence messages automatically route with system monitoring. Uncertain messages queue for prioritized human review. This hybrid approach balances automation with human judgment, maintaining customer experience while protecting channels from spam.

\subsection{Further Developments Necessary for Full Production}

\textbf{Currently Available}:
\begin{itemize}
    \item Trained models ready to use
    \item GUI interface for service teams
    \item Classification engine functional
    \item Documentation complete
\end{itemize}

\textbf{Required for Enterprise Deployment}:
\begin{enumerate}
    \item \textbf{REST API} (Weeks 1-3): Enable system-to-system integration
    \item \textbf{Monitoring Dashboard} (Weeks 4-6): Real-time KPI tracking
    \item \textbf{Automated Retraining} (Weeks 7-9): Periodic model updates
    \item \textbf{Multilingual Support} (Weeks 10+): Extend beyond English
\end{enumerate}

\subsection{Critical Success Factors}

For the spam filter to deliver sustained business value:

\begin{enumerate}
    \item \textbf{Service Team Trust}: Staff must understand and trust model decisions
    \item \textbf{Clear Processes}: Explicit procedures for different confidence levels
    \item \textbf{Feedback Integration}: Regular incorporation of team feedback
    \item \textbf{Realistic Expectations}: Perfection is impossible; trade-offs are inherent
    \item \textbf{Adequate Resources}: Sufficient staff for manual review of uncertain messages
    \item \textbf{Ongoing Monitoring}: Continuous tracking of key metrics
\end{enumerate}

\section{Recommendations}

Deploy the medium-risk configuration for optimal spam protection and customer experience balance. Implement the confidence-based review workflow. Establish daily monitoring discipline. Plan quarterly retraining with accumulated misclassified messages. Move gradually to full system integration as organizational confidence builds. Consider the proposed further developments based on scale and business requirements.

\end{document}
